\section{模型的评价}

\subsection{模型的优点}

\textbf{（1）评价对象与题意一致。} 模型以患者全部项目完整完成为计数单元，未排项目和超过截止的事件均保留在评价分母中。问题二只处理住院患者，问题三再加入门诊和体检需求，两种场景的完成率具有明确的资源边界。

\textbf{（2）设备与准备条件具有现实对应。} 22台设备按题目表1逐项目建立兼容集合，复合项目按组成项目确定共同能力；床旁位置与项目能力取交集。空腹、憋尿、工作时段和跨室转运均落实到具体项目，避免把不同检查路径视为同质任务。分时段、分项目组预约在实际超声流程中也具有改善候检秩序的应用基础\cite{chenjuan2023,nhc2020}。

\textbf{（3）统计参数和调度决策衔接完整。} 需求阈值、项目目录与医生容量由训练期确定，留出期用于评价；题给时长进入主排程，历史报告间隔只辅助判断相对负荷。问题一输出的工时、兼容集合和时隙容量直接构成问题二、三的模型参数，并与超声工作量和设备质量管理指标的业务口径相衔接\cite{nhc2022}。

\textbf{（4）大规模求解保持患者级完整性。} 构造算法枚举患者内部候选序列，优先搜索共同设备，并在全部项目可行后整体提交。该机制与完整事件目标一致，也保留了设备互斥、医生并发、准备和转运等约束。从留出期抽取并额外构造的代表性子问题用于精确对照，最终排程再通过独立计算重建约束。精确对照只覆盖未触发修复的基线分支，真实峰值日审计则只检查修复的自然触发与约束可行性，两类证据的边界彼此分开。

\textbf{（5）多目标结果便于政策选择。} 问题三把背景服务水平写成全评价期下限，以$\varepsilon$约束生成非支配点。医院可根据门诊、体检政策下限选择对应方案，并同时观察住院完成率、等待时间和资源利用率，而不必预先给两个目标设置难以解释的权重。

\subsection{模型的不足}

\textbf{（1）报告完成时刻采用代理。} 附件没有扫查结束、报告书写、审核和提交时刻，模型以最后项目结束近似报告完成。因此，排程完成率适合评价预约与扫查资源，尚不能单独刻画报告环节的真实延迟。

\textbf{（2）背景计划日来自回顾性信息。} 门诊、体检报告日期不是原始预约日期，现有压力场景未包含爽约、迟到和改约行为。背景完成率表示标准班次的计划日承载能力，部署时仍需接入真实预约和到达数据。

\textbf{（3）医生能力只刻画总并发。} 题目给出医生总量，但历史医生与未来33名医生之间没有实名映射。当前模型按星期与5分钟槽估计共享并发容量，尚未描述每名医生的专项资质、班次、请假和设备偏好，专项供给可能因此被高估。

\textbf{（4）部分服务参数仍为区间近似。} 多数项目采用题给时长区间，模型尚未利用同部位多项目可能存在的准备共享、扫查协同和报告共享。憋尿与跨室时间通过敏感性给出边界，但患者个体差异仍未进入主模型。

\textbf{（5）全年求解采用构造式算法。} 全年问题规模较大，主模型采用患者级构造与修复控制计算量。8个代表性子问题只为住院优先基线提供局部最优证据，真实峰值日修复审计也不证明修复分支最优；当高背景下限未被当前候选搜索满足时，不能据此认定数学模型不可行。对于极端高负荷和长程资源竞争，仍需用滚动CP-SAT或局部重优化扩大搜索。

\textbf{（6）公平性尚未成为主目标。} 当前词典序首先提高完整事件数，项目少或设备选择较多的患者可能更容易排入。模型虽保留全部患者，但尚未对科室、严重程度和多项目患者设置最低服务率。推荐结果适用于题给设备、标准日班和当前需求结构；资源数量、资质集合或工作时段改变后需要重新求解。
